\documentclass[11pt,a4paper]{article}
\usepackage[utf8]{inputenc}
\usepackage[T1]{fontenc}
\usepackage[english]{babel}
\usepackage{geometry}
\usepackage{hyperref}
\usepackage{graphicx}
\usepackage{enumitem}
\usepackage{xcolor}
\usepackage{titlesec}
\usepackage{fancyhdr}
\usepackage{longtable}
\usepackage{booktabs}

\geometry{margin=2.5cm}

% Define colors
\definecolor{headergreen}{RGB}{15,102,65}
\definecolor{accentred}{RGB}{201,0,48}

% Title formatting
\titleformat{\section}{\Large\bfseries\color{headergreen}}{\thesection}{1em}{}[\titlerule]
\titleformat{\subsection}{\large\bfseries\color{headergreen}}{\thesubsection}{1em}{}

\title{\textbf{Libraries Boosting Connectivity}\\[0.5em]\large Project Documentation}
\author{Technical Documentation}
\date{\today}

\begin{document}

\maketitle
\tableofcontents
\newpage

% ============================================================================
\section{Project Overview}
% ============================================================================

\textbf{Libraries Boosting Connectivity} is an interactive web application designed to visualize and explore global library connectivity data. The project aims to provide a comprehensive picture of the state of digital infrastructure in libraries worldwide.

\subsection{Purpose}

The application serves as a tool for:
\begin{itemize}
    \item Visualizing library connectivity data on an interactive map
    \item Exploring statistics about internet access, connection types, and perceived quality
    \item Filtering and analyzing data by country and library type
    \item Understanding the challenges libraries face regarding digital connectivity
\end{itemize}

\subsection{Technology Stack}

\begin{longtable}{|p{4cm}|p{10cm}|}
\hline
\textbf{Technology} & \textbf{Purpose} \\
\hline
React 19 & Frontend framework for building the user interface \\
\hline
Vite & Build tool and development server \\
\hline
Mapbox GL & Interactive map rendering and visualization \\
\hline
react-map-gl & React wrapper for Mapbox GL \\
\hline
PapaParse & CSV file parsing for data loading \\
\hline
React Router & Client-side navigation between pages \\
\hline
Tailwind CSS & Utility-first CSS framework for styling \\
\hline
\end{longtable}

% ============================================================================
\section{Application Architecture}
% ============================================================================

The application follows a component-based architecture with clear separation between pages and reusable components.

\subsection{Project Structure}

\begin{verbatim}
frontend/
├── public/           # Static assets and CSV data files
│   ├── arxiu_sortida.csv
│   ├── dades.csv
│   └── img/          # Images and icons
├── src/
│   ├── App.jsx       # Main application component with routing
│   ├── main.jsx      # Application entry point
│   ├── components/   # Reusable UI components
│   ├── pages/        # Page-level components
│   ├── assets/       # Application assets
│   └── styles/       # Global styles
└── docs/             # Documentation
\end{verbatim}

% ============================================================================
\section{Pages}
% ============================================================================

The application consists of three main pages accessible through the navigation.

\subsection{Home Page}

The landing page of the application featuring:
\begin{itemize}
    \item Hero section with project introduction and call-to-action buttons
    \item Animated network visualization representing library connections
    \item Statistical summary boxes showing key metrics:
    \begin{itemize}
        \item Number of connected libraries
        \item Libraries using DSL connections
        \item Percentage of libraries satisfied with their connection
    \end{itemize}
    \item Footer with navigation links
\end{itemize}

The statistics are dynamically calculated from the CSV data and link to the Map page with corresponding filters applied.

\subsection{Map Page}

The core of the application, providing an interactive map visualization:
\begin{itemize}
    \item Full-screen Mapbox map with library markers
    \item Side panel with filters and statistics
    \item Multiple visualization modes (library status, connection type, perceived quality, not connected)
    \item Country and library type filtering
    \item Responsive design with mobile-optimized layout
    \item Legend display for current visualization mode
\end{itemize}

\subsection{About Page}

Information page describing the project's mission, methodology, and team.

% ============================================================================
\section{Components}
% ============================================================================

\subsection{Navigation Components}

\subsubsection{Topbar}
The main navigation header displayed on Home and About pages. Features:
\begin{itemize}
    \item Brand logo linking to home
    \item Navigation links (About, Country Profiles, Explore LBC Map)
    \item Responsive hamburger menu for mobile devices
\end{itemize}

\subsubsection{TopBrand}
Compact brand component used on the Map page. Features:
\begin{itemize}
    \item Hamburger menu with dropdown navigation
    \item Brand logo and title
    \item Menu closes when clicking outside
\end{itemize}

\subsubsection{Footer}
Page footer with:
\begin{itemize}
    \item Partner logos
    \item Navigation links
    \item Copyright information
\end{itemize}

\subsection{Map Components}

\subsubsection{MenuLateral (Side Panel)}
The main control panel for the map, featuring:
\begin{itemize}
    \item Country and library type dropdown filters
    \item Statistics display (InfoCards) showing connectivity metrics
    \item Connection type distribution charts
    \item Perceived quality visualization
    \item Bottom filter buttons for switching visualization modes
    \item Collapsible design with smooth animations
    \item Mobile-responsive bottom sheet layout
\end{itemize}

\subsubsection{MapLegend}
Dynamic legend component that changes based on the active visualization mode:
\begin{itemize}
    \item \textbf{Library Status}: Connected/Not connected/Unknown
    \item \textbf{Type of Connection}: Fiber, DSL, Satellite, Cable, Mobile data, etc.
    \item \textbf{Perceived Quality}: Very poor to Excellent scale
    \item \textbf{Not Connected}: Reasons for lack of connectivity
\end{itemize}

\subsubsection{LibraryInfoPopup}
Popup displayed when hovering over a library marker showing basic library information.

\subsubsection{LibraryDetailsPanel}
Detailed information panel shown when a library is selected. Displays:
\begin{itemize}
    \item Library name, type, and country
    \item Data source information
    \item Expandable sections with detailed connectivity data
    \item Target audience information
    \item Services affected by connectivity issues
\end{itemize}

\subsubsection{LibrarySectionCard}
Collapsible card component used within LibraryDetailsPanel to organize information into expandable sections.

\subsection{Filter Components}

\subsubsection{FilterHeader}
Header section of the side panel containing:
\begin{itemize}
    \item Title and share button
    \item Country dropdown selector
    \item Library type dropdown selector
\end{itemize}

\subsubsection{BottomFilters}
Row of filter buttons for switching between visualization modes:
\begin{itemize}
    \item Library Status (default)
    \item Type of Connection
    \item Not Connected
    \item Perceived Quality
\end{itemize}

Each filter has an info tooltip explaining its purpose.

\subsection{Statistics Components}

\subsubsection{StatBox}
Card component displaying a statistic with:
\begin{itemize}
    \item Large number with optional unit
    \item Description text
    \item Link to explore more on the map
    \item Decorative image
\end{itemize}

\subsubsection{InfoCard}
Compact statistics card used in the side panel showing:
\begin{itemize}
    \item Metric value and description
    \item Icon representation
    \item Distribution bars (for quality metrics)
\end{itemize}

\subsection{Visual Components}

\subsubsection{AnimatedNetwork}
Decorative animated visualization on the Home page representing a network of connected nodes, symbolizing library connectivity.

% ============================================================================
\section{Data Flow}
% ============================================================================

\subsection{Data Sources}

The application loads data from CSV files stored in the \texttt{public/} directory:
\begin{itemize}
    \item \textbf{dades.csv}: Main library data with coordinates and connectivity information
    \item \textbf{arxiu\_sortida.csv}: Processed data for statistics
\end{itemize}

\subsection{Data Processing}

When the application loads:
\begin{enumerate}
    \item CSV files are fetched and parsed using PapaParse
    \item Data is transformed into GeoJSON format for map visualization
    \item Statistics are computed (connection counts, quality percentages, etc.)
    \item Data is categorized into buckets (download speed, connection type, quality)
\end{enumerate}

\subsection{Filtering}

Users can filter data by:
\begin{itemize}
    \item \textbf{Country}: Shows only libraries from selected country
    \item \textbf{Library Type}: Filters by library category (Public, Academic, etc.)
    \item \textbf{Visualization Mode}: Changes how markers are colored on the map
\end{itemize}

% ============================================================================
\section{Visualization Modes}
% ============================================================================

\subsection{Library Status (Default)}

Colors libraries based on internet access:
\begin{itemize}
    \item \textcolor{green}{\textbf{Green}}: Connected to internet
    \item \textcolor{red}{\textbf{Red}}: Not connected
    \item \textcolor{cyan}{\textbf{Cyan}}: Unknown status
\end{itemize}

\subsection{Type of Connection}

Shows the type of internet connection each library uses:
\begin{itemize}
    \item Fiber optic, DSL, Satellite, Cable, Mobile data, Other, Unknown
\end{itemize}

\subsection{Perceived Quality}

Visualizes how libraries rate their digital infrastructure:
\begin{itemize}
    \item Scale from Very Poor (0-19) to Excellent (80-100)
\end{itemize}

\subsection{Not Connected}

Shows only libraries without internet and their reasons:
\begin{itemize}
    \item Infrastructure limitations
    \item High cost
    \item Electrical supply issues
    \item Digital literacy gaps
    \item Policy/Regulatory barriers
\end{itemize}

% ============================================================================
\section{Responsive Design}
% ============================================================================

The application is fully responsive with specific adaptations for mobile devices:

\subsection{Desktop Layout}
\begin{itemize}
    \item Side panel on the left with collapsible toggle
    \item Legend positioned at bottom-right of map
    \item Full filter controls visible
\end{itemize}

\subsection{Mobile Layout}
\begin{itemize}
    \item Bottom sheet panel that can be expanded/collapsed
    \item Floating info button to show legend in modal
    \item Optimized zoom padding to account for bottom panel
    \item Touch-friendly interactions
\end{itemize}

% ============================================================================
\section{URL Parameters}
% ============================================================================

The Map page supports URL parameters for deep linking:

\begin{longtable}{|p{4cm}|p{10cm}|}
\hline
\textbf{Parameter} & \textbf{Description} \\
\hline
\texttt{filter} & Sets the initial visualization mode. Valid values: \texttt{library\_status}, \texttt{type\_connect}, \texttt{not\_connect}, \texttt{perceived\_quality} \\
\hline
\end{longtable}

Example: \texttt{/map?filter=perceived\_quality}

% ============================================================================
\section{User Interactions}
% ============================================================================

\subsection{Map Interactions}

\begin{itemize}
    \item \textbf{Click on library}: Selects library and shows details panel
    \item \textbf{Click on country}: Zooms to country and filters data
    \item \textbf{Click on sea/empty area}: Resets to worldwide view
    \item \textbf{Hover on library}: Shows popup with basic info
\end{itemize}

\subsection{Panel Interactions}

\begin{itemize}
    \item \textbf{Country dropdown}: Filter libraries by country
    \item \textbf{Library type dropdown}: Filter by library category
    \item \textbf{Bottom filter buttons}: Switch visualization mode
    \item \textbf{Arrow toggle}: Expand/collapse side panel (desktop) or bottom sheet (mobile)
\end{itemize}

% ============================================================================
\section{Deployment}
% ============================================================================

The application is configured for deployment on Vercel with the following setup:

\begin{itemize}
    \item Build command: \texttt{npm run build}
    \item Output directory: \texttt{dist}
    \item Environment variables required: \texttt{VITE\_MAPBOX\_TOKEN}
\end{itemize}

% ============================================================================
\section{Future Considerations}
% ============================================================================

Potential areas for expansion:
\begin{itemize}
    \item Country Profiles page (currently linked but not implemented)
    \item Additional data visualization charts
    \item Data export functionality
    \item Multi-language support
    \item Real-time data updates
\end{itemize}

\end{document}
